% GENERATED FILE. Edit canonical lessons, then run npm run generate:books.
% canonical-source-hash: fnv1a64:dd3e245a0c37c882
% canonical-lessons: FR-C518-exposition, FR-C518-vernissage, FR-C518-galerie, FR-C518-oeuvre, FR-C518-encadrer

\chapter{Une exposition}
\label{ch:fr-une-exposition}
\hlchaptermodality{518}

\begin{chapteropening}
\textbf{By the end of this chapter:} \emph{I can talk about going to an exhibition.}

Complete the last lesson of chapter 518: I can talk about going to an exhibition.
\end{chapteropening}

\section[l'exposition]{l'exposition — exhibition}

\label{lesson:FR-C518-exposition}



\begin{quote}
Before the new one: say the French for marble, then the French for sculptor.
\end{quote}

\subsection*{You'll want to know: l'exposition}
\textbf{l'exposition} — \textquotedblleft{}exhibition\textquotedblright{}.

\begin{quote}
Nous avons vu une exposition de photos.
\end{quote}

\begin{quote}
\emph{We saw an exhibition of photos.}
\end{quote}

Feminine. Informally \emph{une expo}.

\begin{cousinweb}
\emph{l'exposition} is a learned borrowing from Latin \emph{expositiō} \textquotedblleft{}a setting out\textquotedblright{} (EtymDB), from \emph{expōnere} \textquotedblleft{}to put out\textquotedblright{}, built on \emph{pōnere} \textquotedblleft{}to put\textquotedblright{} (Lewis \& Short).
\end{cousinweb}

\begin{grammarlens}[title={What you've built}]
One more word for telling, explaining and describing at B1.
\end{grammarlens}

\subsection*{Your turn}
\begin{itemize}
\raggedright
  \item \emph{Say it:} \emph{l'exposition}
  \item \emph{Say it:} \emph{l'exposition}, once more
  \item \emph{Say it:} a sentence of your own with \emph{l'exposition}
  \item \emph{Recall:} say the French for marble, then the French for sculptor, then say \emph{l'exposition} again
\end{itemize}

\begin{tcolorbox}[breakable,colback=teal!4,colframe=teal!35!black,title={Before you move on}]
What does l'exposition mean? (\textquotedblleft{}Exhibition\textquotedblright{}.) Say it once more.
\end{tcolorbox}
\section[le vernissage]{le vernissage — private view, opening (of an exhibition)}

\label{lesson:FR-C518-vernissage}



\begin{quote}
Before the new one: say the French for sculptor, then the French for exhibition.
\end{quote}

\subsection*{You'll want to know: le vernissage}
\textbf{le vernissage} — \textquotedblleft{}private view, opening (of an exhibition)\textquotedblright{}.

\begin{quote}
Elle nous a invités au vernissage de son exposition.
\end{quote}

\begin{quote}
\emph{She invited us to the opening of her exhibition.}
\end{quote}

First the last coat of varnish, put on the day before the show.

\begin{cousinweb}
\emph{le vernissage} is built on \emph{vernir} \textquotedblleft{}to varnish\textquotedblright{} with \emph{-age} (EtymDB): painters once varnished their works on the eve of the show. The origin of \emph{vernis} is usually traced to Medieval Latin \emph{veronix}.
\end{cousinweb}

\begin{grammarlens}[title={What you've built}]
One more word for telling, explaining and describing at B1.
\end{grammarlens}

\subsection*{Your turn}
\begin{itemize}
\raggedright
  \item \emph{Say it:} \emph{le vernissage}
  \item \emph{Say it:} \emph{le vernissage}, once more
  \item \emph{Say it:} a sentence of your own with \emph{le vernissage}
  \item \emph{Recall:} say the French for sculptor, then the French for exhibition, then say \emph{le vernissage} again
\end{itemize}

\begin{tcolorbox}[breakable,colback=teal!4,colframe=teal!35!black,title={Before you move on}]
What does le vernissage mean? (\textquotedblleft{}Private view, opening (of an exhibition)\textquotedblright{}.) Say it once more.
\end{tcolorbox}
\section[la galerie]{la galerie — gallery; arcade}

\label{lesson:FR-C518-galerie}



\begin{quote}
Before the new one: say the French for exhibition, then the French for private view.
\end{quote}

\subsection*{You'll want to know: la galerie}
\textbf{la galerie} — \textquotedblleft{}gallery; arcade\textquotedblright{}.

\begin{quote}
Cette galerie vend les tableaux de jeunes artistes.
\end{quote}

\begin{quote}
\emph{This gallery sells paintings by young artists.}
\end{quote}

\emph{Une galerie marchande} — a shopping arcade.

\begin{cousinweb}
\emph{la galerie} was borrowed from Italian \emph{galleria}, from Medieval Latin \emph{galeria}. The earlier origin is not settled.
\end{cousinweb}

\begin{grammarlens}[title={What you've built}]
One more word for telling, explaining and describing at B1.
\end{grammarlens}

\subsection*{Your turn}
\begin{itemize}
\raggedright
  \item \emph{Say it:} \emph{la galerie}
  \item \emph{Say it:} \emph{la galerie}, once more
  \item \emph{Say it:} a sentence of your own with \emph{la galerie}
  \item \emph{Recall:} say the French for exhibition, then the French for private view, then say \emph{la galerie} again
\end{itemize}

\begin{tcolorbox}[breakable,colback=teal!4,colframe=teal!35!black,title={Before you move on}]
What does la galerie mean? (\textquotedblleft{}Gallery; arcade\textquotedblright{}.) Say it once more.
\end{tcolorbox}
\section[l'œuvre]{l'œuvre — work (of art)}

\label{lesson:FR-C518-oeuvre}



\begin{quote}
Before the new one: say the French for private view, then the French for gallery.
\end{quote}

\subsection*{You'll want to know: l'œuvre}
\textbf{l'œuvre} — \textquotedblleft{}work (of art)\textquotedblright{}.

\begin{quote}
Le musée garde les œuvres d'un grand peintre.
\end{quote}

\begin{quote}
\emph{The museum holds the works of a great painter.}
\end{quote}

Feminine for a single work: \emph{une œuvre d'art}. \emph{Un chef-d'œuvre} — a masterpiece.

\begin{cousinweb}
\emph{l'œuvre} is inherited from Latin \emph{opera} \textquotedblleft{}work, effort\textquotedblright{} (EtymDB), related to \emph{opus} \textquotedblleft{}a work\textquotedblright{} (Lewis \& Short).
\end{cousinweb}

\begin{grammarlens}[title={What you've built}]
One more word for telling, explaining and describing at B1.
\end{grammarlens}

\subsection*{Your turn}
\begin{itemize}
\raggedright
  \item \emph{Say it:} \emph{l'œuvre}
  \item \emph{Say it:} \emph{l'œuvre}, once more
  \item \emph{Say it:} a sentence of your own with \emph{l'œuvre}
  \item \emph{Recall:} say the French for private view, then the French for gallery, then say \emph{l'œuvre} again
\end{itemize}

\begin{tcolorbox}[breakable,colback=teal!4,colframe=teal!35!black,title={Before you move on}]
What does l'œuvre mean? (\textquotedblleft{}Work (of art)\textquotedblright{}.) Say it once more.
\end{tcolorbox}
\section[encadrer]{encadrer — to put in a frame; to supervise}

\label{lesson:FR-C518-encadrer}



\begin{quote}
Before the new one: say the French for gallery, then the French for work.
\end{quote}

\subsection*{You'll want to know: encadrer}
\textbf{encadrer} — \textquotedblleft{}to put in a frame; to supervise\textquotedblright{}.

\begin{quote}
J'ai fait encadrer une vieille photo de mes parents.
\end{quote}

\begin{quote}
\emph{I had an old photo of my parents framed.}
\end{quote}

Also to supervise: \emph{encadrer une équipe}.

\begin{cousinweb}
\emph{encadrer} is built from \emph{en-} and \emph{cadre} \textquotedblleft{}frame\textquotedblright{}, borrowed from Italian \emph{quadro}, from Latin \emph{quadrum} \textquotedblleft{}a square\textquotedblright{} (Lewis \& Short).
\end{cousinweb}

\begin{grammarlens}[title={What you've built}]
One more word for telling, explaining and describing at B1.
\end{grammarlens}

\subsection*{Your turn}
\begin{itemize}
\raggedright
  \item \emph{Say it:} \emph{encadrer}
  \item \emph{Say it:} \emph{encadrer}, once more
  \item \emph{Say it:} a sentence of your own with \emph{encadrer}
  \item \emph{Recall:} say the French for gallery, then the French for work, then say \emph{encadrer} again
\end{itemize}

\begin{tcolorbox}[breakable,colback=teal!4,colframe=teal!35!black,title={Before you move on}]
What does encadrer mean? (\textquotedblleft{}To put in a frame; to supervise\textquotedblright{}.) Say it once more.
\end{tcolorbox}
